\chapter*{关于本手册}
\addcontentsline{toc}{chapter}{关于本手册}
\markboth{关于本手册}{关于本手册}

\section*{这本手册想解决什么问题}

每年九月，都有大量应届本科生在同一个时间窗口里反复问同样的问题：我的排名够不够、夏令营优营到底算不算数、系统里那个待录取通知能在手里放多久、联系导师没回复是不是就没戏了。

这些问题分散在教育部文件、上千个学院的通知、学长学姐的口口相传和各类机构的营销文案里。真正麻烦的是：\textbf{同一个问题，不同学校、不同学院、不同年份的答案是不一样的}，而网络上流传的很多“常识”其实是把某校的局部规则当成了普遍规律。

本手册的做法是：先把真实会被问到的问题尽量完整地拆出来，再回到一手来源去查每一问的答案，然后\textbf{明确标出哪些是全国统一规定、哪些是某校某院的特定规则、哪些只是经验判断}。

\section*{三条阅读约定}

\begin{keybox}
\textbf{一、规则分三层。} 每条问答的正文只写有来源支撑的规则性内容；主观判断、技巧、趋势预判一律放进“说明（经验性，非官方规则）”一行，与规则严格分开。看到这一行，就说明该内容没有官方条文背书。

\textbf{二、来源可回溯。} 正文中的每个 \texttt{[数字引用编号]} 编号都指向书末《来源文献总表》的对应条目，条目里有发布机构、文件类型、标题、年份和原始网址。全书的规则性陈述都可以按这个编号查回原文。

\textbf{三、时效写在明处。} 相关政策结论在正文中标注数字引用，完整文件信息见书末参考文献。推免的政策、系统时间轴、名额和金额每年都可能变化，\textbf{引用旧年份文件的地方会显式标注}。
\end{keybox}

\section*{必须说明的局限}

\begin{itemize}
\item \textbf{招生规则以招生单位当年正式文件为准。} 本手册整理的是公开可查的文件内容，但它不能替代你所在学校教务处、学院和研究生院发布的当年度通知。任何与本手册不一致的地方，以官方文件为准。
\item \textbf{院校规则无法穷举。} 全国有数百所具有推免资格的高校、数以千计的招生院系，本手册只能选取有代表性的样本作为例证。凡引用具体学校的规则，都点明了学校名称与文件年份，\textbf{这些例子不代表其他学校也如此}。
\item \textbf{部分通道确实没有公开依据。} 少数招生类型在教育部层面没有统一公开文件。遇到这类问题，本手册如实写明“未见统一规定”，而不是用推测填补。
\item \textbf{信息截止时间。} 本手册整理的信息截止于 \handbookdate。此后发布的文件、调整的时间节点与金额标准，均需以最新公告为准。
\item \textbf{本手册不提供任何录取承诺。} 也不存在能够影响录取结果的付费服务；凡声称有“内部名额”“保过通道”的，均不可信。
\end{itemize}

\section*{编写方法}

本手册采用“问题空间驱动”的编写方式，而非按传统教材目录组织：

\begin{enumerate}
\item \textbf{拆解问题}：先大规模枚举真实考生的提问，形成问题全集；
\item \textbf{聚类成章}：让章节结构从问题的自然聚类中产生，而不是先定目录再找内容；
\item \textbf{逐条取证}：对每一条问题，从教育部、学信网、省级教育考试机构、高校研究生院与学院官网检索一手材料；
\item \textbf{标注适用边界}：为每条结论标注是全国统一、院校特定还是经验性，并写明依据年份；
\item \textbf{保留冲突}：当不同来源对同一件事说法不一致时，如实并列呈现，而不是挑一个写成定论。
\end{enumerate}

\vspace{1em}
\noindent{\small\color{cMeta}轻保研整理。内容依据公开资料编写，来源信息见书末；政策与院校要求请以最新正式文件为准。}
